Under severe class imbalance, a classifier can achieve high overall accuracy
while performing poorly on the minority class that imbalance-correction
methods are intended to protect. Methods such as resampling, cost-sensitive
reweighting, and TabPFN posterior rescaling are typically evaluated under the
assumption of reliable labels, yet they may be deployed in settings where
labels are imperfect. Considerably less is known about whether these methods
remain beneficial when minority labels are corrupted. This paper addresses
that gap through a controlled audit of seven correction methods across TabPFN
(threshold correction, DistPFN, DistPFN-T, and downsampling) and XGBoost
(cost-sensitive weighting, random oversampling, and SMOTE). Two directed
label-noise mechanisms, majority-to-minority mislabeling and minority loss,
are injected across 26 UCI datasets, four imbalance ratios, five contamination
levels, and 10 seeds. Five metrics proposed in this paper assess correction
performance, while dataset-level Friedman and pairwise Wilcoxon tests assess
statistical significance. Across 112{,}319 evaluations, every correction
underperformed its baseline under some conditions, with harm rates ranging
from 2.4\% to 47.4\%. Excluding the balanced control lowers the maximum harm
rate to 37.8\%. Some corrections became non-beneficial at contamination levels
as low as 5\%, and method differences were significant in both model families
and both noise directions ($p < 10^{-9}$). TabPFN threshold correction was the
most robust, with the lowest harm rate and the highest mean gain in both
directions. Providing TabPFN's rescaling procedure with the true prior did not
consistently improve performance, suggesting that mislabeled context examples
may drive degradation. Resampling methods duplicate mislabeled examples, while
cost-sensitive and rescaling methods inherit the corrupted class ratio.
Scientifically, the study contributes a design that holds training size fixed,
preventing class ratio from being confounded with data volume. Practically,
the findings suggest that imbalance corrections should be evaluated under
realistic label noise before deployment.
